\section{Inside the network: a descriptive look}
\label{sec:network}

Judges and sentence embeddings see the text; the model's own residual stream
sees the computation. This section is exploratory and descriptive: it reports
correlational geometry of the residual stream with judged windows as
observations (clustered by cell; for the pooled correlations we give
cluster-robust intervals from a bootstrap over cells, and the within-condition
values are descriptive), uses a logit lens for one coarse quantity,
and makes no causal intervention. It characterizes what the judged dimensions
co-vary with in the model's state; it does not establish a mechanism. We re-run every finished stream through its generator
(one forward pass with a KV cache, exact token positions) and capture the
residual at 13 layers (every 4th of 48 for Qwen3-30B-A3B; every 3rd for the
36-layer Qwen3-8B and the 40-layer OLMo-2), mean-pooled over 64-token windows
(stride 32) and mean-centered per layer before cosine geometry, plus a logit
lens at each captured layer (the final norm and unembedding applied to the
intermediate state). Three questions.

\paragraph{N1: where does the stream freeze?} Per-layer trajectory
geometry of the window vectors (mean step between consecutive windows,
explored radius) by condition, and the logit-lens \emph{commitment
layer}, the captured layer from which the top-1 token no longer changes.
Bare generation's mean step is a third to a quarter of every interrupted
condition's at every sampled layer; its explored radius is 0.13 at the input layers
against 0.42--0.53, and the gap narrows with depth but never closes (0.24 vs
0.33--0.34 at layer 47, CIs disjoint; Figure~\ref{fig:n1}). The three
interrupted conditions are geometrically alike inside the network; what
separates them is what the judge reads. On the 30B, bare generation's top-1
stabilizes \emph{later} (commitment index 11.12 \ci{11.04}{11.25} of 12 vs
10.88--10.96) while its final distribution is far more confident (final
entropy 0.34 vs 0.42--1.08), a pattern consistent with copying, a late-layer
computation that ends certain; the ordering does not replicate on the 8B, so
we report it as an observation. The 8B and OLMo replicate the geometry (bare
radius 0.14 and 0.36 at layer 0 against 0.46--0.61).

\begin{figure}[t]
\centering
\includegraphics[width=\linewidth]{hidden_h1_geometry_scaffold.png}
\caption{N1: residual-stream trajectory geometry per layer and condition
(Qwen3-30B-A3B, ablation battery): mean step, explored radius, and logit-lens
commitment. Bare generation moves least at every sampled layer, most so at the surface.}
\label{fig:n1}
\end{figure}

\paragraph{N2: which layer's movement predicts judged surprise?} For
every judged window, its novelty at layer $l$ (cosine distance between the
window's mean state and the mean state of everything before it) and its local
step (vs the previous 160 tokens), correlated with judged surprise
(Figure~\ref{fig:n2}). Pooled over conditions, novelty against the past
correlates with surprise at the input layers ($\rho=+0.47$ at layer 0;
cell-bootstrap 95\% CI \ci{+0.22}{+0.67}) and not at the top ($-0.04$,
\ci{-0.26}{+0.17}), but the pooled number is carried by the
bare/interrupted contrast. Within interrupted conditions the sign flips with
depth: local step at layer 0 correlates positively with surprise
(salience-only $+0.46$, scaffold $+0.23$) and at layer 47 negatively
($-0.34$, $-0.14$); novelty against the whole past is negative at depth
(salience-only $-0.53$ at layer 47). Across the 668 judged windows of
battery~2, layer-0 movement predicts surprise within every condition
($\rho$ +0.3 to +0.6), and higher final entropy predicts it everywhere
($\rho$ +0.3 to +0.6): confident is unsurprising. In these data, judged
surprise co-varies with departure at the surface layers and with continuity
at depth: new at the surface, continuous underneath.

\begin{figure}[t]
\centering
\includegraphics[width=\linewidth]{hidden_h2_novelty_scaffold.png}
\caption{N2: Spearman correlation between the judged window's movement
(novelty vs everything before it, left; local step, right) at each layer and
judged surprise, pooled and within condition.}
\label{fig:n2}
\end{figure}

\paragraph{N3: how deep does an interruption reach?} Cosine distance
between the 64 tokens before an injection and the 64 after it (the injected
text skipped), per layer, minus the same at random positions; and the change
in similarity to the premise's own state (Figure~\ref{fig:n3}). Across a
scaffold reseed \emph{with forgetting}, the state after the injection differs
from the state before by +0.09 (layer 4) growing to +0.21 (layer 40) beyond
control, and its similarity to the premise's state rises by +0.09 to +0.20
with depth: forgetting moves the deep state toward the premise's own state,
a re-encounter with the beginning as the network represents it. Across a clock reseed \emph{over preserved
context} the same measures are +0.01--0.03 and $\approx$0 at every
layer, and this is the interruption the judge rewards most. In this reading,
judged surprise and connection do not track deep representational shifts;
they accompany surface departures over a deep state that is left nearly
intact. We take this as a description of what the judged dimensions track,
not as an account of why the arm scores as it does. The dissociation replicates on the 8B (+0.03 to +0.10 vs
+0.01--0.03) and, larger, on OLMo (+0.20 to +0.38 with a return of +0.13 to
+0.18, vs +0.02--0.06 with no return). In battery 2, the depth of the state
change scales with the length of the interrupted segment (period 75:
$\approx$0; 150: +0.01--0.02; 300: +0.02--0.04; 600: +0.04--0.075) and not
with what is injected (the internal counterpart of the frequency result), and
the salience-timed stitch is the one 150-scale injection that reliably moves
the deep state and pulls it toward the premise, while being the arm whose
stream the judge rates lowest: deep movement and judged quality dissociate
here too. Along the stream, similarity to the premise state is U-shaped in
depth ($\approx$0.9 at layer 0, $\approx$0.5 mid-network, rising at the top)
and ordered at the top layer scaffold 0.85 $>$ salience-only 0.79 $>$ clock
reseed 0.67 $>$ bare 0.55: the sentence-embedding ``return to the premise'' of
the scaffold is visible at the top layer and not in the middle of the network.

\begin{figure}[t]
\centering
\includegraphics[width=\linewidth]{hidden_h3_interruption_scaffold.png}
\caption{N3: how deep an interruption reaches (Qwen3-30B-A3B). Left: state
change across an injection per layer, minus random-position control. Right:
change in similarity to the premise's own state. The forgetting reseed moves
the deep state and returns it toward the premise; the clock reseed over
preserved context barely touches it, and is what the judge rewards.}
\label{fig:n3}
\end{figure}
